%%%%%%%%%%%%%%%%%%%%%%%%%%%%%%%%%%%%%%%%%
% Beamer Presentation
% LaTeX Template
% Version 1.0 (10/11/12)
%
% This template has been downloaded from:
% http://www.LaTeXTemplates.com
%
% License:
% CC BY-NC-SA 3.0 (http://creativecommons.org/licenses/by-nc-sa/3.0/)
%
%%%%%%%%%%%%%%%%%%%%%%%%%%%%%%%%%%%%%%%%%

%----------------------------------------------------------------------------------------
%	PACKAGES AND THEMES
%----------------------------------------------------------------------------------------

\documentclass{beamer}

\mode<presentation> {

% The Beamer class comes with a number of default slide themes
% which change the colors and layouts of slides. Below this is a list
% of all the themes, uncomment each in turn to see what they look like.

%\usetheme{default}
%\usetheme{AnnArbor}
%\usetheme{Antibes}
%\usetheme{Bergen}
%\usetheme{Berkeley}
%\usetheme{Berlin}
%\usetheme{Boadilla}
%\usetheme{CambridgeUS}
%\usetheme{Copenhagen}
%\usetheme{Darmstadt}
%\usetheme{Dresden}
%\usetheme{Frankfurt}
%\usetheme{Goettingen}
%\usetheme{Hannover}
%\usetheme{Ilmenau}
%\usetheme{JuanLesPins}
%\usetheme{Luebeck}
\usetheme{Madrid}
%\usetheme{Malmoe}
%\usetheme{Marburg}
%\usetheme{Montpellier}
%\usetheme{PaloAlto}
%\usetheme{Pittsburgh}
%\usetheme{Rochester}
%\usetheme{Singapore}
%\usetheme{Szeged}
%\usetheme{Warsaw}

% As well as themes, the Beamer class has a number of color themes
% for any slide theme. Uncomment each of these in turn to see how it
% changes the colors of your current slide theme.

%\usecolortheme{albatross}
%\usecolortheme{beaver}
%\usecolortheme{beetle}
%\usecolortheme{crane}
%\usecolortheme{dolphin}
%\usecolortheme{dove}
%\usecolortheme{fly}
%\usecolortheme{lily}
%\usecolortheme{orchid}
%\usecolortheme{rose}
%\usecolortheme{seagull}
%\usecolortheme{seahorse}
%\usecolortheme{whale}
%\usecolortheme{wolverine}

%\setbeamertemplate{footline} % To remove the footer line in all slides uncomment this line
%\setbeamertemplate{footline}[page number] % To replace the footer line in all slides with a simple slide count uncomment this line

%\setbeamertemplate{navigation symbols}{} % To remove the navigation symbols from the bottom of all slides uncomment this line
}

\usepackage{graphicx,caption} % Allows including images
\usepackage{booktabs} % Allows the use of \toprule, \midrule and \bottomrule in tables
\usepackage[utf8]{inputenc} %für ß
\usepackage{subfigure}
\usepackage{amsmath} 
\usepackage{comment}
\usepackage[ngerman,english]{babel}

%----------------------------------------------------------------------------------------
%	TITLE PAGE
%----------------------------------------------------------------------------------------

\title[Convolutional Neural Networks]{Bilderkennung mit Convolutional Neural Networks} % The short title appears at the bottom of every slide, the full title is only on the title page

\author{Ivo Tremel und Lennard Nöhren} % Your name
\institute[CHI] % Your institution as it will appear on the bottom of every slide, may be shorthand to save space
{
Leibniz Universität Hannover \\ % Your institution for the title page
\medskip
%\textit{john@smith.com} % Your email address
}
\date{18. Januar 2018} % Date, can be changed to a custom date

\begin{document}

\begin{frame}
\titlepage % Print the title page as the first slide
\end{frame}

\begin{frame}
\frametitle{Overview} % Table of contents slide, comment this block out to remove it
\tableofcontents % Throughout your presentation, if you choose to use \section{} and \subsection{} commands, these will automatically be printed on this slide as an overview of your presentation
\end{frame}

%----------------------------------------------------------------------------------------
%	PRESENTATION SLIDES
%----------------------------------------------------------------------------------------

%------------------------------------------------
\section{Convolutional Neural Network} % Sections can be created in order to organize your presentation into discrete blocks, all sections and subsections are automatically printed in the table of contents as an overview of the talk
%------------------------------------------------
\begin{frame}
	\frametitle{Convolutional Neural Network}
	\begin{itemize}
		\item Besonders für Verarbeitung von Bilddaten geeignet
		\item Lokale Korrelation der Daten bildet Motive
		\item Convolutional Layers entdecken lokale Features
		\item Pooling Layers führen Features zusammen
		\item Parallelitäten zur Hierarchie des visuellen Cortex
	\end{itemize}
\end{frame}

\subsection{Convolutional Layer} % A subsection can be created just before a set of slides with a common theme to further break down your presentation into chunks


\begin{frame}
	\frametitle{Convolutional Layer}
	 \begin{itemize}
		\item Diskrete Faltung auf dem Eingabebild
		\item 5x5 Matrix läuft über das Bild und filtert die Werte
		\item Matrix heißt Filterkernel
		\item Betrachtung der lokalen Umgebung
		\item Features können aus mehreren Punkten bestehen
	\end{itemize}
	
\end{frame}
\begin{frame}
	\frametitle{2D Discrete Convolution}
	\[
	\begin{bmatrix}
	1 & 2 & 3 \\
	4 & 5 & 6  \\
	7& 8 & 9 
	\end{bmatrix} * 
	\begin{bmatrix}
	1 & 2 & 1 \\
	0 & 0 & 0  \\
	-1& -2 & -1 
	\end{bmatrix} = 
	\begin{bmatrix}
	-13 & -20 & -17 \\
	-18 & -24 & -18  \\
	13& 20 & 17 
	\end{bmatrix}
	\]  
	\begin{itemize}
		\item \[ y[m,n]=\sum_{j=-\infty}^\infty\sum_{i=-\infty}^\infty x[i,j]*h[m-i,n-j]\]
	\end{itemize}
	\hspace*{0pt}\hbox{\scriptsize Quelle:\thinspace{\small\itshape http://songho.ca/dsp/convolution/convolution.html\#convolution\_2d}}
\end{frame}


%------------------------------------------------

\subsection{Pooling Layer}


\begin{frame}
	\frametitle{Pooling Layer}
	\begin{itemize}
		\item Fassen den Output von vorhergehenden Neuronen zusammen
		\item Bei MaxPooling wird der größte Wert durchgeschaltet
		\item Andere Werte werden verworfen
		\item Reduktion der Datenmenge ermöglicht schnellere Berechnung
	\end{itemize}
\end{frame}
\begin{frame}
	\frametitle{MaxPooling}
	\begin{figure}
		\includegraphics[width=0.9\linewidth]{images/maxpool.jpeg}
	\end{figure}
	\hspace*{0pt}\hbox{\scriptsize Quelle:\thinspace{\scriptsize\itshape https://github.com/cs231n/cs231n.github.io/blob/master/assets/cnn/maxpool.jpeg}}
\end{frame}


\begin{frame}
	\frametitle{Aufbau}
	\begin{figure}
		\includegraphics[width=0.4\linewidth]{images/main_graph_conv_reduced.png}
		\caption{Aufbau klassisches CNN}
		\label{fig:main_graph_conv}
	\end{figure}
\end{frame}

%------------------------------------------------
\section{Inception Module}
%------------------------------------------------

\begin{frame}
	\frametitle{Inception Module}
	\begin{figure}
		\includegraphics[width=0.8\linewidth]{images/inception_meme.jpg}\\
		\hspace*{0pt}\hbox{\scriptsize Quelle:\thinspace{\small\itshape http://knowyourmeme.com/memes/we-need-to-go-deeper}}
		\label{fig:inception_meme}
	\end{figure}
	
\end{frame}

\begin{frame}
	\frametitle{Inception Module}
	\begin{columns}[c]
		\column{0.5\textwidth}
		\begin{itemize}
			\item Analyse des Inputs in verschiedenen Skalierungen
			\item Für niedrige Level 1x1 Convolution für kleine Features
			\item 3x3 und 5x5 Convolution für größere Features
			\item Max Pooling sehr effektiv in CNN's \\ 
			$\rightarrow$ Als Alternativer Pfad zu Convolutions enthalten
		\end{itemize}
		
		\column{0.5\textwidth}
		\begin{figure}
			\includegraphics[width=\linewidth]{images/inception_naive_graph.png}
			\caption{Aufbau Inception Module}
			\label{fig:naive_inception}
		\end{figure}
	\end{columns}
\end{frame}

\begin{frame}
	\frametitle{Features aus verschiedenen Layers}
	\begin{figure}
		\includegraphics[width=0.4\linewidth]{images/deep_learning.jpg}\\
		\hspace*{0pt}\hbox{\scriptsize Quelle:\thinspace{\small\itshape https://deliveryimages.acm.org/10.1145/2010000/2000787/figs/f2.html}}
		\caption{Features aus verschiedenen Layers eines CNN's}
		\label{fig:deep_learning_features}
	\end{figure}
\end{frame}

\begin{frame}
	\frametitle{Inception Module}
	\begin{columns}[c]
		\column{0.5\textwidth}
		\begin{itemize}
			\item Sehr hoher Rechenaufwand
			\item Dimension steigt stark an
			\item Deshalb: 1x1 Convolutions vor den größeren Modulen
			\item Reduziert die Dimension des Inputs
			\item Ermöglicht mehrere Inception Module in Folge 
		\end{itemize}
		
		\column{0.5\textwidth}
		\begin{figure}
			\includegraphics[width=\linewidth]{images/inception_pro_graph.png}
			\caption{Aufbau verbessertes Inception Module}
			\label{fig:pro_inception}
		\end{figure}
	\end{columns}
\end{frame}

\begin{frame}
	\frametitle{Aufbau}
	\begin{figure}
		\includegraphics[width=0.4\linewidth]{images/main_graph_inception_reduced.png}
		\caption{Aufbau CNN mit Inception Layer}
		\label{fig:main_graph_inception}
	\end{figure}
\end{frame}

%------------------------------------------------
\section{Datensatz}
%------------------------------------------------

\begin{frame}
	\frametitle{Datensatz}
	\begin{itemize}
		\item 1000 Handschriftliche Buchstaben aus 10 Klassen
		\item Buchstaben A M O T U Y V und Umlaute Ä Ö Ü
		\item Einige falsche gelabelte Bilder
	\end{itemize}
	\begin{figure}
		%\includegraphics[width=0.2\textwidth]{hand_images/image_Ae_up_106_point_24.png}
		\subfigure{\includegraphics[width=0.15\textwidth]{hand_images/image_Ae_up_106_point_24.png}}
		\subfigure{\includegraphics[width=0.15\textwidth]{hand_images/image_M_up_9_point_94.png}}
		\subfigure{\includegraphics[width=0.15\textwidth]{hand_images/image_O_up_5_point_14.png}}
		\subfigure{\includegraphics[width=0.15\textwidth]{hand_images/image_T_up_29_point_14.png}}
		\subfigure{\includegraphics[width=0.15\textwidth]{hand_images/image_U_up_31_point_328.png}}
		\subfigure{\includegraphics[width=0.15\textwidth]{hand_images/image_Ae_up_142_point_31.png}}
		\subfigure{\includegraphics[width=0.15\textwidth]{hand_images/image_Oe_up_157_point_17.png}}
		\subfigure{\includegraphics[width=0.15\textwidth]{hand_images/image_Ue_up_132_point_13.png}}
		\subfigure{\includegraphics[width=0.15\textwidth]{hand_images/image_Y_up_46_point_8.png}}
		\subfigure{\includegraphics[width=0.15\textwidth]{hand_images/image_V_up_46_point_314.png}}	
		\caption{Beispiele für Buchstaben aus dem Datensatz}
		\label{fig:charexample}
	\end{figure}
\end{frame}

\begin{frame}
	\frametitle{Falsch gelabelter Buchstabe}
	\begin{itemize}
		\item \glqq image\_Ä\_up\_106\_point\_24.png\grqq
	\end{itemize}
	\begin{figure}
		\includegraphics[width=0.5\textwidth]{hand_images/image_Ae_up_106_point_24.png}
	
		\caption{Beispiel für falsch gelabelten Buchstaben }
		\label{fig:wrongcharexample}
	\end{figure}
\end{frame}

%------------------------------------------------
\section{Ergebnisse}
%------------------------------------------------
\begin{frame}
	\frametitle{Details zu Lernparametern}
	\begin{itemize} 		
		\item 1000 Epochen mit 750 Trainingsbildern
		\item Unterschiedliche Trainingsgeschwindigkeiten
		\item AdamOptimizer
	\end{itemize}
	\begin{table}
		\begin{tabular}{l l l l}
			\toprule
			\textbf{Batchsize} & \textbf{Lernrate} & \textbf{Weights} & \textbf{Biases}\\
			\midrule
			50 &  0.0001 & \( \mathcal{N}(0,0.1) \) & 0.1\\
			
			\bottomrule
		\end{tabular}
		\caption{Lernparameter}
	\end{table}
\end{frame}

\begin{frame}
	\frametitle{Gelernte Kernel Filter}
	\begin{figure}
	\subfigure{\includegraphics[width=0.2\textwidth]{wl25.pdf}}
	\subfigure{\includegraphics[width=0.2\textwidth]{wl26.pdf}}
	\subfigure{\includegraphics[width=0.2\textwidth]{wl27.pdf}}
	\subfigure{\includegraphics[width=0.2\textwidth]{wl28.pdf}}
	\caption{4/32 vom ersten Convolutional Layer gelernte Convolution Kernels der Größe 5x5x3}
	\label{fig:kernel}
	\end{figure}
	\begin{itemize}
		\item Hohe Werte rötlich, niedrig positive gelb
		\item Nahe bei 0 grün, negative Werte blau 
		\item Abstrakte, high-level Features des Convolutional Layers
		\item Aus Bilderkennung bekannte Filter
	\end{itemize}
\end{frame}

\begin{frame}
	\frametitle{Kantenerkennung Sobel Filter}
	\[ G_{x}=
	\begin{bmatrix}
	1 & 2 & 0 & -2 & -1 \\
	4 & 8 & 0 & -8 & -4 \\
	6 & 12 & 0 & -12 & -6 \\
	4 & 8 & 0 & -8 & -4 \\
	1 & 2 & 0 & -2 & -1 
	\end{bmatrix}
	\text{und } G_{y}=
	\begin{bmatrix}
	1 & 4 & 6 & 4 & 1 \\
	2 & 8 & 12 & 8 & 2 \\
	0 & 0 & 0 & 0 & 0 \\
	-2 & -8 & -12 & -8 & -2 \\
	-1 & -4 & -6 & -4 & -1 
	\end{bmatrix}
	\]
	\centering\hspace*{0pt}\hbox{\scriptsize Quelle:\thinspace{\small\itshape Intels IPP Bibliothek}}
	\begin{figure}
		\centering
		\subfigure[$G_x$ detektiert vertikale Kanten]{\includegraphics[width=0.25\textwidth]{xsobel.pdf}}
		\subfigure[$G_y$ detektiert horizontale Kanten]{\includegraphics[width=0.25\textwidth]{ysobel.pdf}}
		\caption{Kantenerkennung durch Sobel-Operatoren}
		\label{fig:sobel}
	\end{figure}
\end{frame}

\begin{frame}
	\frametitle{Kantenerkennung Sobel Filter}
	\begin{figure}
		\centering
		\subfigure{\includegraphics[width=0.25\textwidth]{wl25.pdf}}
		\subfigure{\includegraphics[width=0.25\textwidth]{ysobel.pdf}}
		\caption{Ähnlichkeit zu Sobel-Operatoren}
		\label{fig:sobel}
	\end{figure}
	\begin{itemize}
		\item Linker Filter erinnert horizontal gespiegelten Sobelfilter $G_y$  
		\item Faltung mit dem Eingabebild erzeugt ein Gradientenbild
	\end{itemize}
\end{frame}

\begin{frame}
	\frametitle{Ergebnisse normales CNN}
	\begin{figure}
		\subfigure[Genauigkeit]{\includegraphics[width=0.4\textwidth]{images/1000_step_accuracy_conv_small.png}}
		\subfigure[Cross Entropy]{\includegraphics[width=0.4\textwidth]{images/1000_step_cross_entropy_conv_small.png}}
		\label{fig:conv_results}
	\end{figure}
	\begin{itemize}
		\item Trainingsgenauigkeit nach ungefähr 750 Epochen bei 100\%
		\item Testgenauigkeit nach dem Training ungefähr 85\%
		\item Sehr schnelles Training
	\end{itemize}
\end{frame}

\begin{frame}
	\frametitle{Ergebnisse CNN mit Inception Layer}
	\begin{figure}
		\subfigure[Genauigkeit]{\includegraphics[width=0.4\textwidth]{images/1000_step_accuracy_inception_small.png}}
		\subfigure[Cross Entropy]{\includegraphics[width=0.4\textwidth]{images/1000_step_cross_entropy_inception_small.png}}
		\label{fig:inception_results}
	\end{figure}
	\begin{itemize}
		\item Trainingsverlauf ähnlich wie bei dem anderen Netz
		\item Gelegentliche Einbrüche der Cross Entropy während Training
		\item Testgenauigkeit auch ungefähr bei 80\%-85\%
		\item Training dauert deutlich länger
	\end{itemize}
\end{frame}

%------------------------------------------------
\section{Diskussion}
%------------------------------------------------

\begin{frame}
	\frametitle{Diskussion}
	\begin{itemize}
		\item Niedrige Genauigkeit im Vergleich zu CNN's mit MNIST Datenset
		\item Grund dafür: Größe des Datensets und fehlerhafte Labels
		\item Hier nur sehr kleine Netze verwendet, Größere Netze sind effektiver für Bildanalyse
		\item Kein Mechanismus um Overfitting zu reduzieren
		\item Netz mit Inception Modul bringt keine besseren Ergebnisse
		\item Vorteil Von Inception: in niedrigen Layern kleine Convolutions, in höheren Layern größere Convolution
		\begin{itemize}
			\item In Netz mit nur einem Layer nicht sinnvoll
			\item Vermutlich Nicht notwendig für simple Bilder wie Buchstaben
		\end{itemize}
	\end{itemize}
\end{frame}

%------------------------------------------------
\section{Zusammenfassung und Ausblick}
%------------------------------------------------

\begin{frame}
	\frametitle{Zusammenfassung und Ausblick}
	\begin{itemize}
		\item 85\% Genauigkeit für Schrifterkennung mit kleinem CNN
		\item Inception Module bringen keinen Vorteil in diesem Kontext
		\item Datenset nicht gut für das Training von Neuronalen Netzen geeignet
		\item Für zukünftige Untersuchungen:
		\begin{itemize}
			\item Datenset vergrößern und verbessern
			\item Untersuchen ob größere Netze effizienter sind
			\item Prüfen ob Inception Module in größeren Netzen für Schrifterkennung eine Verbesserung bringen
			\item Mechanismus gegen Overfitting einbauen
		\end{itemize}		
	\end{itemize}
\end{frame}

%------------------------------------------------

\begin{frame}
\Huge{\centerline{Vielen Dank für die}}
\Huge{\centerline{Aufmerksamkeit}}
\end{frame}

%-------------------------------------------------

\end{document} 